% !TEX root = ../main.tex
\section{Related Work}
\label{sec:related-work}

This chapter reviews prior work on the Fischertechnik Industry~4.0 learning factory
(Section~\ref{sec:related-work-fischertechnik}), on the object-to-image (O2I) ratio and
the reliance of convolutional neural networks (CNNs) on the image background
(Section~\ref{sec:related-work-o2i}), on explainability methods that attribute a
decision to an object or to its context (Section~\ref{sec:related-work-explainability}),
on methods that reduce such reliance (Section~\ref{sec:related-work-mitigation}) and on
industrial material classification and small objects
(Section~\ref{sec:related-work-industrial}), before Section~\ref{sec:related-work-gap}
states the gap that remains.

% ============================================================================
\subsection{Prior Work on the Fischertechnik Learning Factory}
\label{sec:related-work-fischertechnik}

S.~Schneidereit et al.~\cite{schneidereitFischertechnik2023} applied YOLO (You Only Look
Once) object detection~\cite{redmon2016yolo} to the learning factory to monitor a
simulated manufacturing process. The learning factory is a model production line with
five modules, namely a sorting machine, an oven, a milling machine, a high-bay warehouse
and a vacuum gripper. The dataset of that study is not the one used in this thesis. It
consists of 700 webcam images in seven batches of 100 that show only white, blue and red
workpieces, with 1,200 labelled workpieces in total. The batches were arranged so that
each colour appears in four of them, to avoid a colour bias in training. Six pixel-wise
augmentations, namely Gaussian blur, brightness, linear contrast, motion blur, noise and
sharpness, enlarged the dataset to 4,900 images, which were split in the proportions
80/10/10. The degrading augmentations account for environmental variation, while the
enhancing ones are meant to compensate for the colour correlation between the workpieces
and the learning factory, which the authors describe as inherent to the task. The
experiments compare four detectors of different sizes, YOLOv3s, YOLOv3l, YOLOv5s and
YOLOv5l, two training and validation strategies and the six augmentations; the abstract
also names different prior-shape assignment strategies. Red workpieces had the lowest
mean average precision at an intersection-over-union (IoU) threshold of 0.5 (mAP@[.5])
in all cases, which the authors attribute to the strong colour correlation with
the learning factory, whose parts are mostly red and blue. They preferred YOLOv3l over
the two YOLOv5 models for detecting the red workpieces, and sharpness was the most
effective augmentation against the colour correlation. In their conclusion, YOLOv3l and
YOLOv5s are each best suited to a particular purpose: YOLOv3l when workpieces share
their colour with the background, and YOLOv5s for strongly distorted images. The earlier
of the two studies thus already links detection errors in the learning factory to a
workpiece sharing its colour with the background, which is the theme the later study and
this thesis take up.

T.~Schneidereit, Gohrenz and Breu{\ss}~\cite{schneidereitYOLOv8} created the dataset
used in this thesis and studied how the complex background of the learning factory
affects the detection of 15 workpieces of different materials and colours. Each
workpiece was photographed at the same 33 positions in each of the five learning factory
models, or stations, which gives 165 images per workpiece and 2,475 in total, labelled
with labelImg in the YOLO format. Between recordings, the lighting was varied and the
camera position was changed, so the workpieces were photographed from different angles
and distances, with the learning factory as background. At a given position, viewpoint
and lighting are the same for all 15 workpieces. Of the 165 images per workpiece, 120
were used for training and 45 as evaluation images. The training images form nested
tiny, small and medium sets of 40, 80 and 120 originals, which 13 augmenters enlarged to
1,600, 3,200 and 4,800 images. On each set, a single-class YOLOv8n (nano) and YOLOv8x
(extra large) model was trained per workpiece, i.e., $15 \times 3 \times 2 = 90$ models,
and two multi-class models were added, 92 in total. Detection accuracy was measured
mainly with mAP@[.5:.95], the mean average precision averaged over IoU thresholds from
0.5 to 0.95 in steps of 0.05, which the authors describe as informative about how well an object is
localised, and heatmaps were computed with
layer-wise relevance propagation (LRP)~\cite{bach2015lrp, montavon2019lrp}.

The authors found that workpieces with similar characteristics can behave quite
differently, and that some background components influenced the detection while others
with the same features did not. On the evaluation images, the YOLOv8n models detected
the two yellow workpieces most accurately, which the authors relate to their colour not occurring in the
background. Black plastic scored worst, most likely because of the many black parts of
the learning factory, and this explanation partially holds for the two red workpieces.
Grey plastic, however, did better than black and red, although the base plate is grey,
so the authors conclude that the relation between workpiece and background colour cannot
be the only factor. Their LRP heatmap for red acrylic glass also marks red elements of
the learning factory. In their conclusion, aluminium, steel, black plastic, red plastic
and red acrylic glass were problematic, while white and yellow plastic, orange acrylic
glass and light wood were learned very well. Besides colour and contour, the reflecting
characteristic of a workpiece is named as relevant. The copy of this dataset used here
is described in Section~\ref{sec:background-dataset}.

Set against the material ranking of this thesis (Tab.~\ref{tab:ovr-full-results},
Section~\ref{sec:results-primary-table}), these findings agree in part, although task
and metric differ. Aluminium and steel, the two metals the 2025 study found problematic,
are the two weakest materials here at every O2I ratio. Brass, which was chosen for its
resemblance to the yellow workpieces, was difficult in neither study. The two red
materials, both problematic there, are among the weaker materials here as well. Two differences remain. Grey plastic, which
did better than black and red there, is also among the weaker materials here, and black
plastic, the hardest workpiece there, is here one of the strongest materials and the most
stable across ratios. Part of the black-plastic difference may come from task and
metric. A whole-scene detector has to find the black workpiece among many black parts
and box it tightly, and mAP@[.5:.95] rewards that localisation, whereas the crops here
are centred on the annotated workpiece and F1 scores only the material decision. The
heatmaps of this thesis come from two methods: occlusion
sensitivity~\cite{zeilerFergus2014}, which the thesis topic prescribes, and
Gradient-weighted Class Activation Mapping (Grad-CAM)~\cite{selvaraju2017gradcam}, a
gradient-based method that reads the whole feature map at once. For black plastic, the
Grad-CAM evidence spreads into the background at 30\% and 10\%, in an example crop into
the dark structure above the workpiece, while the classification stays strong
(Section~\ref{sec:results-qualitative}). The black parts of the learning factory thus
seem salient for a black-plastic model, as the 2025 study proposed, although here they
do not lower the classification. The heatmaps of the 2025 study come from LRP, which is
not compared here, so the explanations of the two studies are not directly comparable.
Section~\ref{sec:discussion-material} returns to the comparison with the per-material
values.

% ============================================================================
\subsection{The Object-to-Image Ratio and Background Reliance in the Wider Literature}
\label{sec:related-work-o2i}

Pawlowski et al.~\cite{pawlowski2019needles} study the classification of tiny objects in
large images, e.g., small lesions in medical images, and treat the O2I ratio as an
experimental variable. They define it by area, as the area of the object divided by the
area of the image, which is also the definition used in this thesis. Their testbed
consists of a digit-based dataset (nMNIST) and a histopathology dataset (nCAMELYON). In
nMNIST, the original $28 \times 28$-pixel digits are kept and the ratio is lowered by
enlarging the canvas from $64 \times 64$ to $1024 \times 1024$ pixels, i.e., from 19.1\%
to 0.075\%, with clutter digits at a roughly constant density. The authors note that
downscaling the digits instead might reduce the quality of the object. In nCAMELYON,
crops of 128, 256 and 512 pixels are cut around real lesions so that the lesion falls
into one of the ratio ranges (100--50)\%, (50--10)\%, (10--1)\% and (1--0.1)\%, with real
tissue as background and a lesion size that varies between crops. The tested CNNs reach
reasonable test accuracy above a ratio of 0.3\% on nMNIST and above 1\% on nCAMELYON.
For a fixed dataset size, there is an O2I limit below which no tested model performs
better than chance, and the amount of training data a model needs to generalise grows
rapidly with the inverse of the ratio. Adapting the receptive field to the known object
size is beneficial, although on nCAMELYON the smallest receptive field performed best,
and at very small ratios the choice of global pooling operation affects the
optimisation. In their supplementary material, the authors also use absolute
input-gradient saliency maps as a localisation check. On nMNIST, the average precision
of locating the digit from these maps falls as the ratio falls, for every pooling
method, which the authors note is in line with the classification accuracy, and for
nCAMELYON the maps are shown qualitatively. They do not compare explainability methods
and do not measure how the saliency is distributed beyond its peak.

Their saliency result is a related precedent for the decline of the Grad-CAM centre
measure found in this thesis as the workpiece shrinks
(Section~\ref{sec:results-divergence}), with two differences. In their case, the
saliency fails where the classifier fails, whereas here the Grad-CAM centre measure
falls while the classification declines only moderately and the occlusion centre
measure does not fall at all. Their measure is also the average precision of the
saliency maximum, not a per-pixel concentration such as the centre measure $c$. The
setups differ as well. In nMNIST the object keeps its pixel size, whereas here the crop
grows around the workpiece and is resized to $256 \times 256$ pixels, so the workpiece
covers fewer pixels as the ratio falls, and the O2I ratio and the scale of the workpiece
change together. The nCAMELYON ranges (100--50)\% and (50--10)\% cover the ratios of this
thesis, and there the models of Pawlowski et al.\ still reached reasonable accuracy. A
fall to chance level is therefore not expected at the ratios studied here, although
their limit depends on the dataset size, and each model here is trained on at most 116
positive crops, all showing the same physical workpiece. The mean test F1 over the 15
materials does decline, but gradually, from 0.820 at 90\% to 0.720 at 10\%
(Section~\ref{sec:results-primary-table}).

Xiao et al.~\cite{xiao2021noise} approach background reliance from the other direction.
Instead of shrinking the object, they keep the foreground and manipulate the background
of ImageNet~\cite{deng2009imagenet} images, separating the two with the ImageNet
bounding-box annotations and, where the shape of the foreground is needed, with GrabCut
segmentation. Models reach non-trivial accuracy from the background alone, and for
87.5\% of foregrounds, some background of another class makes the model predict that
class. Background information can thus outweigh the object in a decision. They also
find that more accurate models both exploit background correlations more and are more
robust to changed backgrounds, a finding at the level of whole models. This thesis looks
at a finer
level. Since a separate model is trained for every material and ratio, 75 models in
total, background reliance can be compared between materials and between ratios. Steel
shows why such a comparison needs more than one method. For steel, occlusion assigns
more importance per pixel to the region outside the reference square than to the square
itself at every ratio. This reading is exposed to a possible confound, since the grey
occlusion patch may resemble the metal and remove little of its evidence. Grad-CAM gives
a different picture: for steel its centre measure rises as the ratio falls, and in an
example crop it concentrates on the bright rim of the workpiece and moves partly to the
background only at 10\% (Section~\ref{sec:results-qualitative}). Whether steel relies
on the background therefore remains open. The evidence rests on ten analysed crops per
model and is indicative only, and Section~\ref{sec:discussion-material} returns to it.

Geirhos et al.~\cite{geirhos2019texture} report a bias of a different kind. CNNs trained
on ImageNet classify images mainly by local texture rather than by global shape, whereas
human perception relies predominantly on shape. Training on Stylized-ImageNet, a
style-transferred variant that removes texture cues while preserving shape, shifts the
networks towards shape and, according to the authors, also improves their robustness to
several types of image distortion. Texture bias concerns the type of feature a network
uses, not where in the image that feature lies, so it can occur even when the evidence
lies on the object. Among the papers reviewed here, it is the closest analogue to the
behaviour of aluminium in this thesis. For aluminium, Grad-CAM favours the workpiece at
every ratio, most strongly at 90\% and 70\%, while occlusion with the grey patch shows
no clear preference for workpiece or background
(Section~\ref{sec:results-qualitative}). Classification is weak at every ratio. A
texture-related reading is possible: a polished surface reflects its surroundings and
may offer an unstable texture even when the workpiece is located correctly. This fits
the conclusion of the 2025 study that the reflecting characteristic of a workpiece
matters~\cite{schneidereitYOLOv8}, which this thesis keeps as a hypothesis for both
aluminium and steel. A simpler explanation competes with it, because aluminium is partly
undertrained at 100 epochs (Section~\ref{sec:results-200epoch}). The analogy is also
loose, since the networks studied by Geirhos et al.\ were trained on ImageNet, whereas
the networks here are trained from scratch on a small industrial dataset.

% ============================================================================
\subsection{Explainability Methods and Context Attribution}
\label{sec:related-work-explainability}

Adhikari et al.~\cite{adhikari2024lostincontext} ask how the context of an object
affects feature attribution methods, not only accuracy. Using segmentation masks to
separate object and context regions in feature maps, they introduce a quantitative
measure of object-versus-context attribution. They report that replacing the context of
an image altogether harms performance more than perturbing the existing context. One
reading of this result is that a network encodes which background it expects with a
given object, beyond how much background is present. In control images without
informative context, attribution methods still assign more than 30\% of their importance
to the uninformative background. They also find that object size alone does not predict
how much a model relies on context. This is consistent with an observation in this
thesis, where the placement of heatmap evidence seems to depend on the material as well
as on the ratio, e.g., for steel compared with black plastic
(Section~\ref{sec:discussion-material}).

Dreyer et al.~\cite{dreyer2023revealing} introduce L-CRP, an extension of Concept
Relevance Propagation (CRP)~\cite{achtibat2023crp} to localisation models, which computes a context score for
each prediction of a segmentation or detection model. The score quantifies how much of
the decision is attributable to background concepts rather than object concepts. With
it, they reveal specific biases, e.g., a segmentation model that uses a hurdle as
context for segmenting a horse. Since CRP extends LRP and L-CRP is also applied to YOLO
detectors, the method is related to the LRP heatmaps of the 2025
study~\cite{schneidereitYOLOv8}. The centre
measure $c$ of this thesis (Eq.~\eqref{eq:center-frac},
Section~\ref{sec:methodology-centerfrac}) is a simpler relative. It also reduces a
heatmap to one number about object versus background, but it compares the mean heatmap
importance per pixel inside and outside a reference square, i.e., spatial regions of the
input, whereas L-CRP works with concepts learned inside the network. Extending the
analysis towards concepts, i.e., asking what the network uses in the background and not
only where its evidence lies, is a possible refinement that is not pursued here.

Ennab and Mcheick~\cite{ennab2025interpretability} compare a pixel-level
interpretability method with Grad-CAM for the detection of COVID-19 in chest
images.\footnote{The paper describes its data both as chest radiographs and as chest
computed tomography (CT) scans.} They report that the two methods differ in granularity.
Grad-CAM gives coarser, region-level explanations, which they find insufficient for
subtle pathological detail, while their pixel-level method localises more accurately
and, in their evaluation, goes together with higher diagnostic accuracy. For this
thesis, the relevant point is only that two attribution methods applied to the same
task disagree, in a domain far from industrial material
classification. Their comparison does not involve occlusion sensitivity, and the share
of background in their images is not varied.

% ============================================================================
\subsection{Mitigating Background and Shortcut Reliance}
\label{sec:related-work-mitigation}

Three papers address how background or contextual reliance can be reduced. This thesis
remains diagnostic and does not attempt such an intervention, but the papers indicate
what could follow it. Singh et al.~\cite{singh2020dontjudge} propose two training-time
interventions against contextual bias. The first minimises the spatial overlap between
the class activation maps~\cite{zhou2016cam} of categories that frequently occur
together. The second splits the learned features into separate subspaces for the
category and for its context. They report accuracy gains of up to 24.2 percentage points on their most
context-biased categories when objects appear outside their usual surroundings. The
first intervention acts on class activation maps, the family of maps to which Grad-CAM
belongs, so a Grad-CAM diagnosis such as the one in this thesis could in principle feed
into such a training objective.

Bassi et al.~\cite{bassi2024isnet} propose ISNet, which optimises LRP heatmaps during
training. Ground-truth segmentation masks are used to penalise relevance that falls
outside the object region. They report that this substantially improves
out-of-distribution generalisation on several medical imaging benchmarks, without added
cost at inference time. ISNet uses LRP, the method behind the heatmaps of the 2025
study~\cite{schneidereitYOLOv8}, and turns an explainability-based diagnosis into a
training objective. For the learning factory dataset, the bounding boxes give only a
box-shaped approximation of the masks that ISNet requires, because a box also contains
some background around the cylindrical workpiece. A background-suppression experiment
along these lines could therefore only be approximated, and
Section~\ref{sec:discussion-future} returns to it.

Le et al.~\cite{le2025shortcut} add a caveat to such efforts. They introduce a neuron
spurious score. For each neuron of the penultimate layer, a Grad-CAM heatmap is
binarised with a threshold, and the share of the highlighted region that overlaps a
mask of the spurious feature is averaged over samples. The score is thus a mask-based
overlap share, a relative of the centre measure $c$, which instead compares mean
importance per pixel. In experiments on the Waterbirds and ISIC datasets with ResNet and
Vision Transformer models pre-trained on ImageNet, Deep Feature
Re-weighting~\cite{kirichenko2023dfr} raises the worst-group accuracy. This is the
accuracy on the group of images for which it is lowest, typically one in which the
spurious feature does not match the class. The distribution of the scores, however,
stays the same. The authors read this as a change in how the
neurons are combined rather than a removal of the spurious features, and they
hypothesise that last-layer re-weighting would fail when a model has under-learned the
relevant information. In this thesis, the caveat may apply mainly to aluminium. Its
Grad-CAM evidence favours the workpiece at every ratio, and it is partly undertrained,
so a method that steers evidence towards the object may not be what it needs. Its weak
classification could reflect a weak object signal, which would be related to the
under-learning that Le et al.\ describe, although there it concerns the representation
and here partly the training length. Steel
would be a candidate for background suppression only if its occlusion result reflects
reliance on the background rather than the grey-fill confound.

% ============================================================================
\subsection{Industrial Material Classification and Small-Object Recognition}
\label{sec:related-work-industrial}

Penumuru et al.~\cite{penumuru2020material} classify materials in an Industry~4.0
setting. Their machine-vision pipeline distinguishes aluminium, copper, medium-density
fibreboard and mild steel from surface images, using RGB colour features and classical
machine learning, namely support vector machines, compared with decision trees, random
forests, logistic regression and $k$-nearest neighbours. They report reliable
performance across different camera angles, lighting conditions and lens settings.
Aluminium and steel also occur among the 15 workpieces of this thesis, where they are
the two hardest materials. The imaging setup of Penumuru et al.\ does not appear to
include a complex background like that of the learning factory, which may explain part
of the contrast. The data of this thesis, however, cannot separate the possible causes
of the metals' difficulty, namely their reflecting surfaces, the complex background and,
for aluminium, the training length. The role of reflections, named in the 2025
study~\cite{schneidereitYOLOv8}, therefore remains a hypothesis, which
Section~\ref{sec:results-qualitative} illustrates on single crops.

Nikouei et al.~\cite{nikouei2025smallobject} survey recent work on small object
detection and name low resolution, occlusion, background interference and class
imbalance as core challenges of the field. Some of the definitions they collect are
relative: an object that covers less than 1\% of the image area counts as small. By this
criterion, even the 10\% crops of this thesis are far from small objects. In the
original $960 \times 720$-pixel photographs, however, a box of median size,
$50 \times 46$ pixels, covers about 0.33\% of the image. For a typical workpiece, the
whole-scene detection task of the 2025 study~\cite{schneidereitYOLOv8} is therefore a
small-object task by this criterion, although the largest boxes exceed 1\%.
That study itself names small objects in front of a complicated background as a known
difficulty in its review of related work. The thesis
studies the transition from crops that consist mostly of workpiece to crops that
consist mostly of background, not the detection of tiny objects. Two points of the
survey nevertheless apply. The downsampling inside a CNN removes discriminative detail
such as textures and edges, which may matter for workpieces that differ mainly in their
surface. Small objects in complex backgrounds, industrial ones included, also suffer from
a low signal-to-noise ratio, the framing that Pawlowski et al.\ use as well. Class
imbalance in the survey means that small objects are underrepresented in datasets. The
imbalance in this thesis is of a different kind. It follows from the one-vs-rest
design, in which each binary model sees about 14 negative crops for every positive one,
and it is handled by weighting the loss by class and by selecting models by F1
(Section~\ref{sec:methodology-training}). Background interference, in turn, is the
subject of the explainability pipeline (Section~\ref{sec:methodology-explainability}).

% ============================================================================
\subsection{The Gap This Thesis Fills}
\label{sec:related-work-gap}

The reviewed papers cover the elements of this thesis separately. Pawlowski et
al.~\cite{pawlowski2019needles} vary the O2I ratio systematically, on synthetic digit
canvases and on real tissue, and use one gradient-based saliency method as a
localisation check that degrades with the ratio, but they compare no explainability
methods and do not measure how the evidence is distributed beyond the saliency peak. The two studies on the
learning factory~\cite{schneidereitFischertechnik2023, schneidereitYOLOv8} work with
its complex industrial background, but for detection and, in the 2025 study, with LRP
as the only heatmap method. Xiao et al.~\cite{xiao2021noise} and Geirhos et
al.~\cite{geirhos2019texture} demonstrate background and texture biases through
manipulations of the images, namely background swaps and stylisation, which are
experiments rather than explainability methods. Adhikari et
al.~\cite{adhikari2024lostincontext} measure how attribution is divided between object
and context, but they change the context of an image rather than the share of the image
that the object covers. Ennab and Mcheick~\cite{ennab2025interpretability} compare two
explainability methods directly, but in images where the share of background is not
varied.

What none of these papers does is to compare two explainability methods on the same
crops as a function of a controlled O2I ratio, against a complex industrial background
and with a per-pixel measure of whether the evidence concentrates on the object or on
the background. In this
thesis, occlusion sensitivity and Grad-CAM are applied to the same validation crops of
each of the 75 models. The two methods diverge as the workpiece shrinks, most clearly at
10\%, where the Grad-CAM centre measure has fallen while that of occlusion has not
(Section~\ref{sec:results-divergence}).
% agreement at 90% depends on the occlusion normalisation, decided 20 Oct
A possible mechanism is that occluding a single patch cannot reveal reliance on a large,
repetitive background, because the rest of that background stays visible. This
explanation is offered as a hypothesis, not as a measured cause.
